\section{Introduction}\label{sec: Intro}

\subsection{Main results}

This paper studies the hitting-time mixing problem for the random \(k\)-cycle walk on \(\mathfrak S_n\). Starting from the identity, at each step we multiply by a uniformly chosen \(k\)-cycle, and we stop the walk at the first time \(\tau\) when every card has been touched. Our main result shows that, in the regime
$$
2\leq k=o\left(\frac{n}{(\log n)^4}\right),
$$
this natural stopping time is already a mixing time: \(X_\tau\) is close in total variation to the relevant uniform measure, subject only to the necessary parity constraint. This extends the hitting-time mixing theorem of Jain-Sawhney \cite{jain2024hitting} for random transpositions to random \(k\)-cycles in a range where \(k\) is allowed to grow with \(n\).

The starting point for this problem is the classical theory of mixing for random walks generated by conjugacy classes of \(\mathfrak S_n\). For random transpositions, Diaconis-Shahshahani \cite{diaconis1981generating} introduced the nonabelian Fourier-analytic method and proved sharp cutoff at time \((n\log n)/2\). Subsequent works extended and refined this circle of ideas for more general conjugacy-class walks, including character bounds, cutoff, and limiting profiles for cycle shuffles; see, for example, Roichman \cite{Roichman1996}, Lulov-Pak \cite{lulov2002rapidly}, Hough \cite{hough2016random}, Teyssier \cite{teyssier2020limit}, and Nestoridi-Olesker-Taylor \cite{nestoridi2022limit}. A more refined question asks whether the walk is already mixed at the first time when every card has been touched. This is the earliest possible time at which total-variation mixing can occur, since before this time at least one card has never been moved. For random transpositions, this hitting-time mixing problem was raised by Berestycki and formulated explicitly by Teyssier \cite[Conjecture 1.2]{teyssier2020limit}; it was recently proved by Jain-Sawhney \cite{jain2024hitting}. Their argument combines a precise fixed-time approximation of the walk with a strong-stationary-time type coupling after the cutoff window, and they noted that the method should extend beyond transpositions, in particular to random \(k\)-cycles for fixed \(k\). The present paper carries this out for \(k\)-cycles, and proves hitting-time mixing throughout the growing range above.

Let us now define the random $k$-cycle walk precisely. Let $P_n$ be the
probability measure on $\mathfrak S_n$ which is uniform on the conjugacy class
of $k$-cycles:
\begin{equation}\label{eq: measure of k-cycle}
P_n(\sigma)=
\begin{cases}
\frac{k}{n(n-1)\cdots(n-k+1)}, & \text{if } \sigma \text{ is a } k\text{-cycle},\\
0, & \text{otherwise}.
\end{cases}
\end{equation}
We consider the discrete-time Markov chain
$$
X_{t+1}=\sigma_{t+1}\cdot X_t,
$$
where $\sigma_1,\sigma_2,\ldots$ are independent random variables with law
$P_n$. Thus $X_t\sim P_n^{*t}$.

Because a $k$-cycle has sign $(-1)^{k-1}$, the limiting uniform measure has
to be interpreted with the appropriate parity constraint. Following the set-up
used by Nestoridi and Olesker-Taylor \cite{nestoridi2022limit} and Hough
\cite{hough2016random}, we compare the walk with the uniform measure on the
parity class on which it is supported: if $k$ is odd, then the walk is
supported on $\mathfrak A_n$; if $k$ is even, then the walk is supported on
$\mathfrak A_n$ at even times and on $\mathfrak A_n^c$ at odd times. We
denote by $U_{\mathfrak S_n}$ the uniform measure on $\mathfrak S_n$, and by
$U_{\mathfrak A_n}$ and $U_{\mathfrak A_n^c}$ the uniform measures on
$\mathfrak A_n$ and $\mathfrak A_n^c$, respectively.

The stopping time of interest is the first time at which every card has been
touched by the walk. More explicitly, write the $l$-th chosen $k$-cycle as
$(i_{l,1}\,\cdots\, i_{l,k})$. We define
\begin{equation}\label{eq: hitting time}
\tau
:=
\min\left\{
t\ge 1:
\min_{1\le j\le n}
\sum_{l=1}^t
\mathbf 1\{j\in\{i_{l,1},\ldots,i_{l,k}\}\}
\ge 1
\right\}.
\end{equation}
Before time $\tau$, at least one card has not been touched, so the walk cannot
be close in total variation to the relevant uniform measure. Our main theorem
shows that this obstruction is the only one in the regime
$2\le k=o(n/(\log n)^4)$: throughout this range, which goes well beyond the
fixed $k$ setting suggested by Jain-Sawhney, the walk is already mixed at the
hitting time. 

\begin{thm}\label{thm: hitting time mixing_main thm}
Suppose $n=\omega(1),2\le k\le o(n/(\log n)^{4})$. With notation as above, we have 
$$d_{TV}(X_\tau,U_{\A_n})\le \exp(-(\log\log n)^{1/2+o(1)})$$
when $k$ is odd, and
$$d_{TV}(X_\tau,U_{\mathfrak{S}_n})\le\exp(-(\log\log n)^{1/2+o(1)})$$
when $k$ is even. 
\end{thm}

The phenomenon captured by \Cref{thm: hitting time mixing_main thm} is part of a broader theme in the study of Markov-chain mixing: in many natural chains, the cutoff time is governed by the disappearance of an explicit local obstruction. This viewpoint goes back at least to the theory of strong stationary times of Aldous-Diaconis \cite{aldous1986shuffling} and Diaconis-Fill \cite{diaconis1990strong}, where one often proves mixing by identifying a stopping time at which all initially visible information has been randomized. Our random $k$-cycle hitting time problem fits naturally into this circle of ideas.

We now give an overview of the proof of \Cref{thm: hitting time mixing_main thm}. As in Jain-Sawhney's treatment of
random transpositions, we begin by proving a fixed-time approximation for the
walk near the cutoff window. The guiding heuristic is that, at time
$\frac{n}{k}(\log n+c)$, \footnote{Throughout the paper, time is always integer-valued. Whenever an expression for time is not an integer, we implicitly take its integer part and suppress the floor notation for notational convenience.} the only visible obstruction to uniformity is the set of cards which have not yet been touched. This set should be approximately Poisson, and conditional on this set, the remaining cards should be almost uniformly shuffled on the appropriate parity class. The following definition introduces the auxiliary measure $\nu_t$ which makes this approximation precise.

\begin{defi}\label{defi: construction of nu}
Suppose we have $t\in\Z$. Set $t':=t-\frac{n\log n}{k}$, and $\gamma_{n,t}:=e^{-kt'/n}$. Let $\nu_t$ denote the measure on $\mathfrak{S}_n$ defined by the following sampling process: first, sample $M_t\in\{0,1,\ldots,n\}$ according to the distribution $\mathbf{P}[M_t=x]=\mathbf{P}(\Pois(\gamma_{n,t})=x)/\mathbf{P}(\Pois(\gamma_{n,t})\le n)$, then sample a uniformly random subset $S_t$ of size $M_t$ in $[n]$. Finally,
\begin{enumerate}
\item When $k$ is odd, sample a uniformly random element of $\mathfrak{A}_{[n]\backslash S_t}$ and view it as an element of $\mathfrak{A}_n$ by fixing all of the elements in $S_t$.
\item When $k$ is even, and $t$ is odd, sample a uniformly random element of $\mathfrak{A}_{[n]\backslash S_t}^c$ and view it as an element of $\mathfrak{A}_n^c$ by fixing all of the elements in $S_t$.
\item When $k$ is even, and $t$ is even, sample a uniformly random element of $\mathfrak{A}_{[n]\backslash S_t}$ and view it as an element of $\mathfrak{A}_n$ by fixing all of the elements in $S_t$.
\end{enumerate}
\end{defi}

The measure $\nu_t$ is designed to isolate the contribution of the untouched
cards. It first chooses the number of untouched cards according to the expected
Poisson law at time $t$, then chooses their locations uniformly, and finally
puts a uniform permutation on the remaining cards, subject only to the necessary
parity constraint. Thus $\nu_t$ is an explicitly tractable proxy for the law of
$X_t$ near the cutoff window.

Our main fixed-time input is that this proxy indeed approximates the random
$k$-cycle walk in total variation. The following theorem refines the second author's result \cite[Theorem 1.6]{shen2026k}: by restricting to the smaller range of $k$ considered here, we obtain the stronger error estimate required for the hitting-time argument.

\begin{thm}[\Cref{thm: restate of X and nu} in the text]\label{thm: X and nu}
Suppose that \(2\le k=o(n/(\log n)^4)\), and that
$t' := t-\frac{n\log n}{k}$ satisfies
$|t'|\le \frac{n}{2k}\log\log\log n$. Then, we have
$$d_{TV}(X_t,\nu_t)=o\left((\log n)^{-5/2}\right).$$
\end{thm}

There is currently no known characterization of a strong stationary time for
$k$-cycle shuffling on $n$ cards. In particular, the time at which all
cards have been touched is not itself a strong stationary time. Nevertheless,
one expects that at this time the distribution of the deck is already close to
uniform, and that, as $n\to\infty$, the distribution should converge to the
uniform distribution in total variation distance. Our next goal is to strengthen
the preceding theorem into a hitting-time mixing result. To achieve this, we introduce a marking scheme starting from a time
\begin{equation}\label{eq: starting time of marking scheme}
    t^*= \frac{n\log n}{k}
    -\frac{n}{2k}\log\log\log n+\frac{4n}{k}\log\log\log\log n .
\end{equation}
By the choice of $t^*$, the set of cards not yet touched by time $t^*$ is small with high probability. We mark all cards that have already been touched, and then continue the process with a carefully chosen rule that marks the remaining
cards one by one. The marking rule is designed so that the time needed to mark all cards is close to the time needed to touch all cards.

The passage from fixed-time to hitting time follows the philosophy introduced by Jain-Sawhney for the random transposition walk. Their key insight is that a sufficiently sharp fixed-time approximation can be bootstrapped into an
``approximate sufficient statistic'' statement: near the cutoff window, once one conditions on the set of untouched cards, the remaining permutation is already close to uniform. Starting from this conditional uniformity, one can then run a
short marking procedure which mimics a strong stationary time only for the few cards that remain unmarked.

We follow this strategy, but the random $k$-cycle walk introduces new difficulties. First, a single update may involve many unmarked cards, so the
marking rule must distinguish the useful event that exactly one new card is absorbed from the exceptional event that several unmarked cards appear in the same $k$-cycle. Second, the parity of a $k$-cycle depends on $k$. When $k$ is
odd, the walk stays in $\A_n$, while when $k$ is even it alternates between $\A_n$ and $\A_n^c$. Thus the artificial marking procedure cannot simply insert idle moves, as in the transposition case; it must preserve the correct parity at
each time.

The main novelty is a parity-compatible marking scheme which resolves these two issues simultaneously. Heuristically, the scheme adds artificial one-card marking moves so that each remaining unmarked card is discovered at the correct rate, while keeping the conditional law on the marked cards uniform on the appropriate parity class. In the even-$k$ case this requires parity-correcting transpositions supported entirely on the marked set. These transpositions change
the parity without affecting which unmarked card is being marked. The proof then compares this artificial process with the genuine $k$-cycle walk and shows that the accumulated discrepancy before all cards are marked is negligible.

\subsection{Future directions.}

The range $2\leq k=o(n/(\log n)^4)$ in
\Cref{thm: hitting time mixing_main thm} should not be interpreted as the
conjectural endpoint of hitting-time mixing for random $k$-cycles. Rather,
this is the regime in which our fixed-time approximation theorem is strong
enough to be combined with the Jain-Sawhney hitting-time argument. We expect
that the same qualitative phenomenon should persist for much larger values of
$k$: once every card has been touched, there should be no further obstruction
to mixing, apart from the necessary parity constraint. The following theorem provides
evidence for this expectation in the opposite, large-cycle regime.

\begin{thm}\label{thm: large cycle}
Suppose $n=\omega(1)$, $k\leq n-1$, and $k\geq n-o(n^{1/2})$. Then, we have
$$
d_{TV}(X_\tau,U_{\mathfrak A_n})=o(1).
$$
\end{thm}

Indeed, as we will see in \Cref{sec: large cycles}, \Cref{thm: large cycle} corresponds to the situation where two cycles
are enough to complete the hitting with high probability. After the first
$k$-cycle, only $n-k=o(n^{1/2})$ cards remain untouched, and a second
independent $k$-cycle covers all of them with probability tending to one. Thus
\Cref{thm: large cycle} shows that, in this regime, the walk is already mixed at the moment
when the second cycle completes the hitting of all cards.

The endpoint $k=n$, however, is degenerate and must be excluded. In that case
the walk hits all cards after a single step, but the distribution at that time
is supported on the conjugacy class of $n$-cycles, and hence is certainly not
close to uniform. With this exception in mind, it is natural to formulate the
following conjecture.

\begin{conj}\label{conj: hitting time mixing all k}
Suppose $n=\omega(1)$ and $2\leq k\leq n-1$. Then the following hold. If $k$ is odd, then
$$
d_{TV}(X_\tau,U_{\mathfrak A_n})=o(1).
$$
If $k$ is even, then 
$$
d_{TV}\!\left(X_\tau,\,\mathbf P(\tau \text{ is even})\,U_{\mathfrak A_n}+\mathbf P(\tau \text{ is odd})\,U_{\mathfrak A_n^c}
\right)=o(1).
$$
\end{conj}

Let us explain the form of the conjecture in the even-\(k\) case. In fact, in \Cref{thm: hitting time mixing_main thm}, when \(k\) is even, we first prove the stronger intermediate statement that \(X_\tau\) is close to the parity mixture
$$
\mathbf P(\tau \text{ is even})\,U_{\mathfrak A_n}+
\mathbf P(\tau \text{ is odd})\,U_{\mathfrak A_n^c}.
$$
We then show, in the regime \(k=o(n/(\log n)^4)\), that the hitting time is asymptotically equally likely to be even or odd, which allows us to replace this mixture by \(U_{\mathfrak S_n}\). For general \(k\), however, the assertion that \(\tau\) has asymptotically balanced parity is no longer expected to hold. Thus the natural even-\(k\) conjecture is convergence to the parity mixture determined by the actual distribution of \(\tau\), rather than necessarily to \(U_{\mathfrak S_n}\).

It would also be interesting to study hitting-time mixing for random walks
generated by more general class measures on $\mathfrak S_n$. In the present
paper the driving measure is supported on a single conjugacy class, namely the
class of $k$-cycles, but the Fourier-analytic framework applies equally well
to any class measure. The main difficulty is then to understand which
irreducible representations make the dominant contribution near the hitting
time, and how the touching process interacts with the representation-theoretic
profile of the chosen class measure. Recent work of Olesker-Taylor, Teyssier,
and Th\'evenin \cite{olesker2025sharp} develops sharp character bounds
and cutoff results for broad families of conjugacy-class walks on
$\mathfrak S_n$, and provides useful representation-theoretic background for
such a direction.

%\JC{Add other directions that you guys come up with here}

\subsection{Notation}

We write \([n]:=\{1,\ldots,n\}\). For a subset \(S\subseteq[n]\), we write \(\mathfrak S_S\) for the group of permutations of \(S\). In particular, when \(S=[n]\setminus T\), we also write \(\mathfrak S_{[n]\setminus T}\) and regard it naturally as a subgroup of \(\mathfrak S_n\) by fixing every point of \(T\). The same convention applies to \(\mathfrak A_S\) and \(\mathfrak A_S^c\).

For a probability measure \(\mu\), we write \(X\sim\mu\) to mean that \(X\) has law \(\mu\), and we write \(\mathcal L(X)\) for the law of a random variable \(X\). The total variation distance between two probability measures \(\mu\) and \(\nu\) on a finite set is denoted by
$$
d_{\mathrm{TV}}(\mu,\nu)
:=
\frac12\sum_x |\mu(x)-\nu(x)|.
$$
When no confusion can arise, we also write \(d_{\mathrm{TV}}(X,Y)\) for \(d_{\mathrm{TV}}(\mathcal L(X),\mathcal L(Y))\). We use \(\mathbf P\) and $\E$ for probability and expectation.

For a permutation \(\sigma\in\mathfrak S_n\), we write
$$
\operatorname{Fix}(\sigma):=\{i\in[n]:\sigma(i)=i\}
$$
for its set of fixed points. If \(a\) is a nonnegative integer and \(b\geq 0\), we use the falling factorial notation
$$
(a)_b:=a(a-1)\cdots(a-b+1),
$$
with the convention \((a)_0=1\).

We write \(f=O(g)\) if \(|f|\le C|g|\) for some absolute constant \(C>0\), \(f=o(g)\) if \(f/g\to0\), and \(f=\omega(g)\) if \(f/g\to+\infty\). All asymptotic notation is understood in absolute value. Moreover, the implicit constants are taken along the given sequence of positive integer inputs; in particular, they are not intended uniformly over all such inputs, but only along the sequence under consideration.

\subsection{Outline of the paper} 

The remainder of the paper is organized as follows. In \Cref{sec: Jain-Sawhney}, we prove the fixed-time approximation theorem, \Cref{thm: X and nu}. In \Cref{sec: fixed to hitting}, we use this approximation to prove the hitting-time mixing result, \Cref{thm: hitting time mixing_main thm}. Finally, in \Cref{sec: large cycles}, we prove \Cref{thm: large cycle} for the regime where \(k\) is close to \(n\).
